\documentclass[12pt,a4paper]{article}
\usepackage{fontspec}
\setmainfont{Times New Roman}
\usepackage[a4paper,margin=2.54cm]{geometry}
\usepackage{setspace}
\onehalfspacing
\usepackage{booktabs,longtable,array,ragged2e,enumitem,titlesec,needspace}
\setlength{\parindent}{0pt}
\setlength{\parskip}{0pt}
\setlist[enumerate]{leftmargin=1.6em,itemsep=0pt,parsep=0pt,topsep=0pt}
\titleformat{\section}{\normalfont\bfseries\fontsize{12}{18}\selectfont}{}{0pt}{}
\titleformat{\subsection}{\normalfont\bfseries\fontsize{12}{18}\selectfont}{}{0pt}{}
\titlespacing*{\section}{0pt}{0.4em}{0.1em}
\titlespacing*{\subsection}{0pt}{0.3em}{0.1em}
\emergencystretch=4em
\begin{document}
\begin{titlepage}
\centering
\vspace*{24mm}
{\bfseries Introduction to Software Engineering Assessment\par}
\vspace{16mm}
{\bfseries CS Group 12 Turnnote Project Proposal\par}
\vspace{8mm}
{\bfseries
Turnnote - A Cross-Platform AI Meeting Assistant for Focused Participation and Actionable Notes\par}
\vspace{18mm}
\begin{tabular}{@{}ll@{}}
\textbf{Group Number:} & CS Group 12 \\
\textbf{Programme of Study:} & BSc (CS) \\
\end{tabular}
\vspace{17mm}
\begin{tabular}{@{}p{.43\textwidth}p{.27\textwidth}@{}}
\toprule
\textbf{Group Member} & \textbf{UoA Student ID} \\
\midrule
Junxuan Zeng & 50106157 \\
\midrule
Rui Lin & 50106129 \\
\midrule
Zitao Hong & 50106160 \\
\midrule
Chenghao Lin & 50106158 \\
\midrule
Zhongtian Lv & 50106136 \\
\midrule
Tianqi Wang & 50106166 \\
\midrule
Jintao Zhang & 50106133 \\
\midrule
Weiqiao Lin & 50106154 \\
\midrule
Dikai Zhang & 50106165 \\
\specialrule{0.8pt}{0pt}{0pt}
Chenxi Wang & 50106132 \\
\bottomrule
\end{tabular}
\vfill
\end{titlepage}
\setcounter{page}{1}
\Needspace{3\baselineskip}
\section*{Problem}
Online meetings are a regular part of professional collaboration. Participants exchange ideas, make decisions, and assign follow-up work, often while trying to maintain their own written record of the discussion. For individual professionals and members of small teams, retaining this information is important for both immediate follow-up and later reference.\par
However, taking detailed notes while participating in a conversation divides attention. A participant who focuses on typing may miss the context of a decision, while someone who concentrates on the discussion may forget important details afterward. Audio recordings preserve the conversation but require additional review, and raw transcripts can be too lengthy to serve as practical meeting notes.\par
The principal problems that Turnnote aims to address are:\par
\begin{enumerate}
\item Divided attention between \textbf{active participation and manual note-taking}.
\item \textbf{Additional effort} required to turn recordings or transcripts into concise summaries and follow-up actions.
\item Difficulty reviewing the source discussion when checking a meeting summary or an assigned task.
\item Fragmented workflows for users working \textbf{across macOS and Windows}.
\item \textbf{Limited visibility} into where meeting content is processed, how it is retained, and how it can be deleted.
\end{enumerate}
These problems motivate a desktop application that brings meeting capture, transcription, note generation, and review into one workflow. The initial target users will be professionals and small teams that participate in recurring online meetings. This target segment and the frequency of the proposed problems will be evaluated through user testing; the proposal does not claim that market demand or productivity improvements have already been measured.\par
\Needspace{3\baselineskip}
\section*{Solution}
We plan to develop \textbf{Turnnote}, a desktop meeting assistant inspired by Granola's product category. Turnnote will support macOS and Windows and use the DeepSeek API to transform meeting transcripts into editable notes. To address the problems above, the application will provide the following functions:\par
\begin{enumerate}
\item \textbf{User-Controlled Meeting Capture:} Users will select their microphone and system-audio sources and explicitly start and stop recording. A visible indicator will show when capture is active, and the application will remind users to obtain any required participant consent before recording.
\item \textbf{Local Speech Transcription:} Turnnote will use a Python \textit{faster-whisper} sidecar process to transcribe meeting audio locally into timestamped English or Chinese text. Raw audio will not be sent to DeepSeek.
\item \textbf{AI Meeting Notes:} On the user's explicit request, the DeepSeek API will organize transcript text into a summary, key topics, decisions, open questions, and action items. Owners and due dates will be included only when supported by the source discussion.
\item \textbf{Review, Editing, and Export:} Users will be able to inspect the transcript, correct generated notes, and export the reviewed result as Markdown.
\item \textbf{Local Meeting Management:} A SQLite database will store meeting records on the user's device. Users will be able to search, rename, and delete records. Audio will be deleted after successful transcription by default, with per-session opt-in retention.
\item \textbf{Cross-Platform Operation and Credential Protection:} A shared desktop interface will provide a consistent workflow on macOS and Windows. DeepSeek API credentials supplied by the user will be stored in the operating system's secure credential store, and the application will explain which text is sent to DeepSeek.
\end{enumerate}
Turnnote will use Electron, React, TypeScript, and Node.js for the desktop application. A Python \textit{faster-whisper} sidecar will provide local speech recognition, while platform-specific audio capture will be accessed through macOS ScreenCaptureKit and Windows WASAPI Loopback integrations. The Electron main process will start and monitor the Python process, manage SQLite, call the DeepSeek API, and coordinate local files and task states. The main process and sidecar will communicate through a controlled JSON Lines interface over standard input and output. The renderer will use a preload bridge with \textit{contextIsolation} enabled, \textit{sandbox} enabled, and \textit{nodeIntegration} disabled. The first release will focus on a complete individual meeting workflow. User accounts, team workspaces, cloud synchronization, calendar integration, automatic meeting detection, meeting bots, speaker diarization, live captions, and external task-management integrations are outside the MVP scope.\par
\Needspace{3\baselineskip}
\section*{Objectives}
\Needspace{3\baselineskip}
\subsection*{Objective 1: Implement meeting audio capture and local transcription, and build a basic cross-platform desktop interface}
\textbf{Core workflow.} The first objective is to establish the core meeting-recording workflow. Users will be able to open Turnnote, select an audio source, record a meeting, and obtain a timestamped transcript that can be reviewed within the application.\par
\textbf{Audio capture.} The first requirement is to implement reliable audio capture on both target platforms. We plan to use ScreenCaptureKit for macOS system audio and WASAPI loopback for Windows system audio, alongside microphone capture. Users will explicitly start and stop recording, and the recording state will remain visible throughout the session. Permission denial, unavailable devices, interrupted capture, and insufficient disk space will produce clear errors and recovery options. Before recording starts, the application will remind the user to obtain any consent required for the meeting.\par
\textbf{Local transcription.} The second requirement is to implement local speech recognition using a Python \textit{faster-whisper} sidecar. The Electron main process will start the sidecar, pass it captured audio through a controlled local interface, receive structured timestamped English or Chinese transcripts, and terminate or restart the process when required. Keeping speech recognition local limits the need to transfer raw meeting audio to an external service, although processing speed and resource use will depend on the user's hardware and selected model. Failed transcription jobs will retain the recoverable meeting input for retry. After successful transcription, temporary audio will be deleted unless the user has opted to retain it.\par
\textbf{Desktop foundation.} The third requirement is to build the desktop application foundation and basic interface. Electron will provide the desktop application shell, React and TypeScript will support the renderer interface, and the Node.js Electron main process will handle application services, IPC, platform integration, and task orchestration. The renderer will not access files, credentials, or network services directly; it will use a controlled preload and IPC bridge. The basic workspace will include audio-source selection, recording controls, processing status, and a transcript view. SQLite will store meeting metadata, transcript segments, and processing state locally.\par
\textbf{Platform baseline.} The planned platform baseline is macOS 13 Ventura or later and Windows 10 version 22H2 or Windows 11, with x64 or arm64 builds where supported by the selected transcription runtime. Platform and hardware combinations will be verified through the release test matrix before being advertised as supported. Successful completion of this objective requires an installable application on both target platforms and a working capture-to-transcript flow with visible recording status and recoverable failure handling.\par
\Needspace{3\baselineskip}
\subsection*{Objective 2: Develop a user-friendly meeting workspace, integrate DeepSeek note generation, and improve processing reliability}
\textbf{Meeting workspace.} The second objective is to turn the basic transcript workflow into a practical meeting workspace. Users will be able to request an AI-generated note, check it against the transcript, and edit the result before using it for follow-up work.\par
\textbf{Interface design.} For the user interface, we will develop the initial UI/UX prototype into a meeting workspace and settings view. The workspace will display the transcript and generated note together so that users can compare them without switching applications. Settings will provide audio-device selection, transcription-model configuration, and DeepSeek API credential management. Recording, transcription, and generation status will remain distinguishable so that users can understand which operation is in progress or requires attention.\par
\textbf{DeepSeek integration.} For AI note generation, we will integrate the DeepSeek API through the Node.js Electron main process. The user will initiate each generation request, and the application will explain that transcript text and the context required for note generation are sent to DeepSeek. Raw audio will remain outside this API flow. User-supplied credentials will be stored in the operating system's secure credential store and excluded from ordinary configuration exports and application logs.\par
\textbf{Structured output.} The generated note will follow a defined structure containing a summary, key topics, decisions, open questions, and action items. An action item's owner or due date will be marked as unspecified when absent from the transcript. Responses will undergo schema validation before being accepted by the application. Schema validation checks the output structure rather than factual correctness, so the transcript and editing tools will remain available for human review.\par
\textbf{Failure handling.} For processing reliability, we will handle timeouts, rate limits, malformed responses, and model context limits explicitly. Long transcripts will be divided into bounded segments and their intermediate results combined into a final note. Retries will be bounded, and a failed request will preserve the local transcript and any existing user edits. Completion of this objective requires successful generation of a structured note from a saved transcript, secure credential handling, and an editing workflow that remains usable when an external request fails.\par
\Needspace{3\baselineskip}
\subsection*{Objective 3: Enhance meeting management and export, and complete cross-platform testing and release preparation}
\textbf{Meeting library.} The third objective is to complete the features needed to manage meeting records over time and prepare Turnnote for pilot use. This will include historical meeting access, export, deletion controls, and verification of the complete application on macOS and Windows.\par
\textbf{Data management.} For meeting management, we will provide a local meeting library with search, renaming, and deletion. Users will be able to reopen transcripts and edited notes and determine whether audio has been retained. Deleting a meeting will remove its application-managed record and associated retained audio. The data-handling notice will distinguish local storage from transcript text transmitted to DeepSeek; deleting a local record does not itself establish deletion from an external provider's systems.\par
\textbf{Markdown export.} For note export, we will provide Markdown output containing the user's reviewed meeting note. Exported content will match the version shown in the editor, including any corrected decisions or action items. This will allow users to transfer meeting results into their existing documentation workflow without requiring a team account or a separate synchronization service.\par
\textbf{Integration testing.} For integration and reliability testing, we will verify capture permissions, long sessions, sleep and resume, network interruption during note generation, local deletion, and application restart. Tests will cover capture, transcription, generation, editing, export, and record management as one end-to-end workflow. API and transcription failures must not silently corrupt or discard saved transcripts or edited notes, and release-blocking defects must be resolved before distribution.\par
\textbf{Pilot evaluation.} For pilot evaluation, we will measure first-workflow completion, note usefulness, processing time, transcription and API failures, repeat use, and processing cost per meeting hour. Initial targets are that at least 80\% of pilot users complete the end-to-end workflow without assistance, average note usefulness reaches 4 out of 5, and at least 95\% of valid DeepSeek requests succeed when the network and service are available. These are evaluation targets rather than achieved results. English and Chinese test meetings will establish the transcription-quality and resource-use baselines.\par
\textbf{Deliverables.} The final deliverables will include macOS and Windows installation packages, a user guide, technical documentation, a privacy and data-handling notice, a known-issues list, and a proof-of-concept evaluation report. Together, these deliverables will demonstrate the completed workflow and document the practical limits of the MVP.\par
\Needspace{3\baselineskip}
\section*{Benefits}
\Needspace{3\baselineskip}
\subsection*{1. More Focused Meeting Participation}
\textbf{User-controlled capture and automated transcription} are intended to reduce the need to type a complete record while taking part in a discussion. A visible recording state gives participants using the application a clear indication of when capture is active.\par
\Needspace{3\baselineskip}
\subsection*{2. More Efficient Post-Meeting Review}
\textbf{Structured notes} will organize lengthy transcripts into summaries, decisions, open questions, and action items. This is intended to reduce the effort required to prepare a useful meeting record, with actual improvements evaluated during the pilot.\par
\Needspace{3\baselineskip}
\subsection*{3. Reviewable and Editable Meeting Results}
Users will be able to compare \textbf{AI-generated content with the transcript} and correct the note before export. This supports human oversight of the final record and helps users identify missing context or unsupported model output.\par
\Needspace{3\baselineskip}
\subsection*{4. Consistent Cross-Platform Workflows}
Turnnote will provide a shared \textbf{meeting workflow on macOS and Windows}. A common interface and Markdown export will make the resulting notes usable across existing documentation practices without requiring all participants to adopt a new collaboration platform.\par
\Needspace{3\baselineskip}
\subsection*{5. Greater Control over Meeting Data}
\textbf{Local transcription, local record storage, optional audio retention, and accessible deletion controls} will give users direct control over application-managed meeting data. Explicit note-generation requests and secure credential storage will make the external AI processing step visible and deliberate.\par
Through these capabilities, Turnnote aims to provide a practical connection between meeting participation, accurate review, and follow-up work. Its value will be assessed against the pilot targets and the reliability of the complete workflow.\par
\Needspace{3\baselineskip}
\section*{Timeline}
\smallskip
\begingroup
\setlength{\tabcolsep}{2pt}
\begin{longtable}{@{}>{\RaggedRight\hyphenpenalty=10000\arraybackslash}p{.12\textwidth}>{\RaggedRight\hyphenpenalty=10000\arraybackslash}p{.59\textwidth}>{\RaggedRight\hyphenpenalty=10000\arraybackslash}p{.24\textwidth}@{}}
\toprule
\textbf{Phase} & \textbf{Description} & \textbf{End Date} \\
\midrule
\endfirsthead
\toprule
\textbf{Phase} & \textbf{Description} & \textbf{End Date} \\
\midrule
\endhead
\bottomrule
\endfoot
Phase 1 & Complete requirements, application architecture, and project initialization & 24 September 2026 \\
\midrule
Phase 2 & Investigate platform audio capture and validate local transcription & 26 September 2026 \\
\midrule
Phase 3 & Develop the Electron desktop interface, Node.js services, and Python transcription sidecar & 20 October 2026 \\
\midrule
Phase 4 & Test interface components, platform audio capture, local transcription, and API handling separately & 10 November 2026 \\
\midrule
Phase 5 & Complete end-to-end integration, cross-platform testing, and pilot evaluation & 20 November 2026 \\
\midrule
Phase 6 & Package and release the MVP, complete documentation, and present the project & 26 November 2026 \\
\end{longtable}
\endgroup
\Needspace{3\baselineskip}
\section*{Action plan}
\Needspace{7\baselineskip}
\subsection*{Objective 1: Implement meeting audio capture and local transcription, and build a basic cross-platform desktop interface}
\smallskip
\begingroup
\setlength{\tabcolsep}{2pt}
\begin{longtable}{@{}>{\RaggedRight\hyphenpenalty=10000\arraybackslash}p{.43\textwidth}>{\RaggedRight\hyphenpenalty=10000\arraybackslash}p{.19\textwidth}>{\RaggedRight\hyphenpenalty=10000\arraybackslash}p{.21\textwidth}>{\RaggedRight\hyphenpenalty=10000\arraybackslash}p{.13\textwidth}@{}}
\toprule
\textbf{Action} & \textbf{Assigned to} & \textbf{Deadline} & \textbf{Progress} \\
\midrule
\endfirsthead
\toprule
\textbf{Action} & \textbf{Assigned to} & \textbf{Deadline} & \textbf{Progress} \\
\midrule
\endhead
\bottomrule
\endfoot
Requirement analysis and requirements documentation & Junxuan Zeng & 22 September 2026 & Completed \\
\midrule
Research and validate macOS ScreenCaptureKit and Windows WASAPI capture & Junxuan Zeng & 22 September 2026 & Completed \\
\midrule
Evaluate English and Chinese transcription with the faster-whisper sidecar & Chenxi Wang & 23 September 2026 & In-progress \\
\midrule
Design the meeting data model and SQLite database structure & Junxuan Zeng & 23 September 2026 & In-progress \\
\midrule
Design and document application interfaces and processing states & Zitao Hong & 23 September 2026 & In-progress \\
\midrule
Initialize the Electron, React, TypeScript, and Node.js application & Weiqiao Lin & 24 September 2026 & In-progress \\
\midrule
Implement capture controls, permissions, timestamped transcription, and the basic workspace & Junxuan Zeng & 24 September 2026 & In-progress \\
\end{longtable}
\endgroup
\Needspace{7\baselineskip}
\subsection*{Objective 2: Develop a user-friendly meeting workspace, integrate DeepSeek note generation, and improve processing reliability}
\smallskip
\begingroup
\setlength{\tabcolsep}{2pt}
\begin{longtable}{@{}>{\RaggedRight\hyphenpenalty=10000\arraybackslash}p{.43\textwidth}>{\RaggedRight\hyphenpenalty=10000\arraybackslash}p{.19\textwidth}>{\RaggedRight\hyphenpenalty=10000\arraybackslash}p{.21\textwidth}>{\RaggedRight\hyphenpenalty=10000\arraybackslash}p{.13\textwidth}@{}}
\toprule
\textbf{Action} & \textbf{Assigned to} & \textbf{Deadline} & \textbf{Progress} \\
\midrule
\endfirsthead
\toprule
\textbf{Action} & \textbf{Assigned to} & \textbf{Deadline} & \textbf{Progress} \\
\midrule
\endhead
\bottomrule
\endfoot
Develop the UI/UX prototype and meeting workspace interface & Chenxi Wang & 25 September 2026 & In-progress \\
\midrule
Define the meeting-note schema and generation prompt & Chenghao Lin & 25 September 2026 & In-progress \\
\midrule
Integrate DeepSeek requests, response validation, and secure credential storage in the Electron main process & Junxuan Zeng, Weiqiao Lin & 26 September 2026 & In-progress \\
\midrule
Implement transcript review and meeting-note editing & Zitao Hong & 27 September 2026 & In-progress \\
\midrule
Implement long-transcript segmentation and summary combination & Chenghao Lin & 27 September 2026 & In-progress \\
\midrule
Test processing performance, bounded retries, and failure recovery & Junxuan Zeng, Rui Lin & 30 September 2026 & In-progress \\
\end{longtable}
\endgroup
\clearpage
\Needspace{7\baselineskip}
\subsection*{Objective 3: Enhance meeting management and export, and complete cross-platform testing and release preparation}
\smallskip
\begingroup
\setlength{\tabcolsep}{2pt}
\begin{longtable}{@{}>{\RaggedRight\hyphenpenalty=10000\arraybackslash}p{.43\textwidth}>{\RaggedRight\hyphenpenalty=10000\arraybackslash}p{.19\textwidth}>{\RaggedRight\hyphenpenalty=10000\arraybackslash}p{.21\textwidth}>{\RaggedRight\hyphenpenalty=10000\arraybackslash}p{.13\textwidth}@{}}
\toprule
\textbf{Action} & \textbf{Assigned to} & \textbf{Deadline} & \textbf{Progress} \\
\midrule
\endfirsthead
\toprule
\textbf{Action} & \textbf{Assigned to} & \textbf{Deadline} & \textbf{Progress} \\
\midrule
\endhead
\bottomrule
\endfoot
Write and structure the project proposal & Junxuan Zeng, Chenghao Lin, Chenxi Wang & 22 September 2026 & Completed \\
\midrule
Implement meeting search, renaming, deletion, and audio-retention controls & Jintao Zhang & 10 October 2026 & In-progress \\
\midrule
Implement Markdown export of reviewed notes & Tianqi Wang & 13 October 2026 & In-progress \\
\midrule
Integrate components and execute macOS and Windows end-to-end tests & Rui Lin, Dikai Zhang & 15 October 2026 & In-progress \\
\midrule
Review credentials, logs, and privacy disclosures & Zhongtian Lv, Dikai Zhang & 20 October 2026 & In-progress \\
\midrule
Write technical documentation, the user guide, and the data-handling notice & Jintao Zhang, Zhongtian Lv, Rui Lin, Chenghao Lin & 10 November 2026 & In-progress \\
\midrule
Conduct the pilot and write the proof-of-concept evaluation report & Tianqi Wang, Zitao Hong, Jintao Zhang & 20 November 2026 & In-progress \\
\midrule
Prepare installation packages, release notes, and the final presentation & Weiqiao Lin, Chenxi Wang, Junxuan Zeng & 26 November 2026 & In-progress \\
\end{longtable}
\endgroup
\end{document}
